% Imporditud: tools/import_eksperimendid.py
% Allikas: Eesti füüsikaolümpiaadi eksperimendi ülesannete
%   kogu (Rael Kalda, 2023), ülesanne Ü15, lahendus L15.
\setAuthor{EFO žürii}
\setRound{piirkonnavoor}
\setYear{2007}
\setNumber{P E1}
\setDifficulty{1}
\setTopic{staatika}
\setVahendid{30 cm pikkune mõõtejoonlaud, 1-kroonine münt (m = 5,0 g)}
\setPunktid{10}
\setAllikas{Eksperimendi kogu Ü15, lahendus lk 38}
\setLahendusseis{toores}

\prob{Joonlaua mass}
Määrata joonlaua mass.

\hint

\solu
Eeldame, et joonlaua mass on kogu joonlaua ulatuses jaotunud ühtlaselt. Paigutame joonlaua risti laua servaga üks ots üle serva. Nihutame joonlauda ettevaatlikult ja määrame koha, kus joonlaud hakkab laua serval kalduma. Sellel kaugusel asub joonlaua masskese. Kuna mõõtejoonlaua skaala ei pruugi alata ega lõppeda joonlaua otstes, ei pruugi masskeskme asukoht olla 15 cm joonel. Mõõdame mündi läbimõõdu (Eesti Panga andmetel 23,25 mm). Asetame mündi mingisse skaala punkti, paigutame joonlaua laua servale ja nihutame joonlauda seni, kuni joonlaud hakkab üle laua serva kalduma. Fikseerime skaalal selle punkti. Arvutame mündipoolse õla dm ja masskeskmepoolse õla dk. Teades mündi massi mm, arvutame seosest Fmdm = Fjdk joonlaua massi mj:

mj = mmdm. dk

Täpsema tulemuse saamiseks tuleb teha mitu katset asetades mündi erinevatele kaugustele masskeskme asukohast ühele ja teisele poole masskeset.

Hindamine: Idee --- (3 p). Joonlaua masskeskme asukoha määramine --- (1 p). Mündi läbimõõdu mõõtmine --- (1 p). Õlgade määramine mõõtmisel müdiga --- (1 p). Kangi seose

teadmine ja kasutamine --- (1 p). Korduvmõõtmised erinevate jõuõlgadega --- (1 p). Korduvmõõtmised juhul kui münt on ühel või teisel pool masskeset --- (1 p). Tulemus (kui erinevus kaalumisel saadud tulemusest ei ületa 5\%) --- (1 p).
\probend
